\begin{frame}{확인 문제 (1)-(2)}
  \small
  \begin{itemize}
    \item \textbf{1.} 시드 0의 학습 도중 마지막 10이 $-673.5$(베이스라인
      $-1274$의 2배)인데, 결정적 평가는 $-1357.9$(베이스라인 \emph{아래})다.
      어느 쪽이 ``배운 정책의 진짜 성능''이고, \textbf{주 지표}는 무엇이어야
      하는가?
      \vskip0.2em
      {\footnotesize $\to$ \textbf{결정적($a=\mu$) 평가}가 탐험을 뺀
      본래 실력(8.1의 $\varepsilon=0$ greedy 평가에 해당). 학습 도중 리턴은
      ``탐험 포함 성과''. 주 지표는 \textbf{결정적 평가} — 평가의 목표가
      ``정책의 안정적인 성능''이지 ``학습 과정이 잘 봤는가''가 아니기
      때문. 학습 곡선은 ``학습이 실제로 되고 있었는가''의 증거로,
      \textbf{역할을 구분해 함께} 보고.}
    \item \textbf{2.} ``연속 정책에는 $\varepsilon$이 없으니, 평가할 때
      $\sigma$ 파라미터를 0으로 설정해서 돌리면 되지 않나요?'' — 왜
      오답인가?
      \vskip0.2em
      {\footnotesize $\to$ $\sigma$는 \textbf{학습된 파라미터}(엔트로피
      보너스와 함께 학습 대상) — $\sigma=0$로 바꾸면 ``원래 정책에서 탐험을
      제거한 실행''이 아니라 \textbf{파라미터를 바꾼 \emph{다른} 정책}의
      실행. 정답: 파라미터를 그대로 두고 \textbf{$a=\mu_\theta(s)$(평균)만
      내보내는 것} — ``같은 정책을 샘플링 없이 실행''하는 것이
      $\varepsilon=0$과 동일한 역할.}
  \end{itemize}
\end{frame}
